\def\version{1}
\problemname{PO Archiving}
After every competition, the Programming Olympiad archives the participants' solutions for all eternity.
However, most solutions submitted during a competition are quite similar.
A contestant who fixes a bug might only change a single line in a solution.
In these cases, it is unnecessary to save both the original solution and the modified solution.
Instead, we can save one of the solutions and the changes that were made between the two solutions.
This can also happen in multiple steps, so that we save a solution $A$, then save the changes between $A$ and another solution $B$, and finally the changes between $B$ and a third solution $C$.

In total, $N$ solutions were submitted during a competition, with sizes $S[0], S[1], \dots, S[N-1]$ bytes.
If the $i$th solution is either saved or can be reconstructed, then the $j$th solution can be reconstructed if the changes that are $D[i][j]$ bytes in size are saved.

In most test groups, the diff sizes will be symmetric, i.e.\ $D[i][j] = D[j][i]$ (similar to regular Unix \texttt{diff} files).
For the last group, however, we consider more general diffs where this does not need to hold.
For example, we can imagine that the diff between the strings \texttt{abaabbaaa} and \texttt{aaaaaa} is stored as ``remove all \texttt{b}'' in one direction, and ``insert \texttt{b} at positions 2, 5, 6'' in the other, where the latter change requires more space to store than the first.

What is the minimum amount of data (in bytes) that must be saved so that all solutions can be reconstructed?

\section*{Input}
The first line contains the integer $1 \le N \le 100$.
The next line contains the $N$ integers $1 \le S[0], S[1], \dots, S[N-1] \le 1\,000\,000$.
The next $N$ lines contain $N$ integers each.
The $i$th of these lines contains the values $1 \le D[i][0], D[i][1], \dots, D[i][N-1] \le 1\,000\,000$.
$D[i][i] = 0$ for all $i$.

\section*{Output}
Output a single number --- the minimum amount of data that must be saved in bytes.

\section*{Scoring}
Your solution will be tested on a set of test groups. To earn points for a group, you must pass all test cases in that group.

\noindent
\begin{tabular}{| l | l | l | l |}
\hline
Group & Points & Constraints \\ \hline
1     & 11         & All values $D[i][j]$ (except when $i = j$) are the same \\ \hline
2     & 27         & All values $S[i]$ are the same, $D[i][j] = D[j][i]$\\ \hline
3     & 48         & $D[i][j] = D[j][i]$ \\ \hline
4     & 14         & No additional constraints. \\ \hline
\end{tabular}
